\documentclass{tutodoc}

\usepackage{../preamble.cfg}


\begin{document}

\section{What does that mean in \tdocquote{English}?}

The macro \tdocmacro{tdocinEN} and its starred version are useless for English speakers because they have the following effects.


\begin{tdoclatex}
Cool and top stand for \tdocinEN*{cool} and \tdocinEN{top}.
\end{tdoclatex}


The macro \tdocmacro{tdocinEN} and its starred version are based on \tdocmacro{tdocquote} : for example, \tdocquote{semantic} is obtained via \tdoclatexin|tdocquote{semantic}|\,.


\begin{tdocnote}
    As the text \tdocquote{in English} is translated into the language detected by \thisproj, the macro \tdocmacro{tdocinEN} and its starred version become useful for non-English speakers.
\end{tdocnote}


\begin{tdocimp}	
	For documents written in a language other than English, if one of the packages \tdocpack{babel} and \tdocpack{polyglossia} is loaded, then, in \tdocquote{the quoted content}\,, the English linguistic rules are respected.
\end{tdocimp}

\end{document}
